\section*{Chapitre \ref{ch:fm:crisis-management} --- Gestion de crise}

\begin{solution}{exo:fm:crisis-management:1}
$0.6\times33=\$19.8$ millions bloqués~; $0.4\times33=\$13.2$ millions reviennent à l'heure 48.
\end{solution}

\begin{solution}{exo:fm:crisis-management:2}
Voir la \cref{def:fm:crisis-management:trapped}~: collatéral déposé au titre de la marge, marge excédentaire non encore restituée, titres réhypothéqués, et soldes
qui auraient dû être ségrégués mais ne l'ont pas été.
\end{solution}

\begin{solution}{exo:fm:crisis-management:3}
Heure 23~: sans la défaillance, seules les hausses progressives de marge liées au choc drainent la trésorerie, et rien ne doit être remargé immédiatement.
\end{solution}

\begin{solution}{exo:fm:crisis-management:4}
Les fonds monétaires et le repo procurent des heures à meilleur marché, si bien que la règle gloutonne les retient en premier~; la ligne reste tirée bien avant son
retrait. Si elle ne pouvait pas être tirée du tout, le plan devrait recourir à une seconde coupe de 10~\% et coûterait 3.12~millions de dollars au lieu de 1.42~millions~: une procédure d'approbation
qui défaille en pleine crise vaut ici 1.7~millions de dollars.
\end{solution}

\begin{solution}{exo:fm:crisis-management:5}
Voir la démonstration de la \cref{prop:fm:crisis-management:monotone}. Si la réduction des pertes réalisées était réglée en espèces, une coupe pourrait diminuer la trésorerie dans les heures qui précèdent
celles où elle réduit les appels, et la proposition serait fausse.
\end{solution}

\begin{solution}{exo:fm:crisis-management:6}
Lesquelles de leurs positions existaient encore et où, quelle marge ils avaient déposée et dans quels comptes, et si elle était ségréguée~: le
relevé quotidien du courtier, les registres propres du desk sur les positions et les garanties, et les registres de la chambre de compensation sur les comptes clients du courtier.
\end{solution}

\begin{solution}{exo:fm:crisis-management:7}
Heure 24, avec une trésorerie à $-\$2.8$ millions à son plus bas~: le second cycle de compensation quotidien l'épuise encore.
\end{solution}

\begin{solution}{exo:fm:crisis-management:8}
Les lignes peuvent être retirées précisément dans la crise qui en a besoin, la trésorerie peut être bloquée chez un courtier défaillant, et la marge peut augmenter chez chaque contrepartie à
la fois~; les appels de l'exercice atteignent 78.8~millions de dollars, et l'\omterm{def:fm:treasury-and-funding:horizon}{horizon de survie} dépend du calendrier, non des totaux.
\end{solution}

\begin{solution}{pb:fm:crisis-management:1}
\begin{enumerate}
\item Voir la \cref{def:fm:crisis-management:plan}.
\item Les décisions de gestion qui entourent le correctif technique~: reprendre ou non, que liquider, qui prévenir, et ce que l'incident a coûté.
\item Voir la \cref{def:fm:crisis-management:drill}.
\item Voir la \cref{def:fm:crisis-management:trapped}.
\item Voir l'\cref{dat:fm:crisis-management:record}.
\item Elle a acheté temporairement des gilts à longue échéance, du 28 septembre au 14 octobre, pour rétablir des conditions de marché ordonnées, après que son comité de politique financière
a relevé les risques liés au dysfonctionnement du marché.
\item Ségrégation des actifs des clients~; elle ne protège que si le courtier la respecte et si son régulateur peut le vérifier.
\item Gestion du risque de liquidité pour les appels de marge et de collatéral, tests de résistance de ces appels, et \omterm{def:fm:operations:collateral}{gestion des garanties} avec suffisamment de trésorerie et d'actifs liquides.
\item Une hausse de 60~\% de la \omterm{def:fm:treasury-and-funding:house}{marge maison} le premier jour, la marge de compensation augmentant de 30~\% puis de 50~\% aux cycles quotidiens, le courtier B défaillant à l'heure 6 avec 60~\%
de son collatéral bloqué, le reste restitué à l'heure 48.
\item Heure 6~; heure 23 sans la défaillance.
\item Fonds monétaires (heure 9, 0.02~million de dollars), repo (heure 15, 0.12~million de dollars cumulés), ligne (heure 24, 0.42~million de dollars), une coupe de 10~\% (au-delà de 72 heures,
1.42~millions de dollars).
\item Voir la \cref{prop:fm:crisis-management:monotone}.
\item Les heures ne sont pas additives et les délais contraignent l'ordre~; une recherche exhaustive sur de petits ensembles d'actions est exacte.
\item Approbations impossibles à obtenir, heures limites manquées, systèmes en panne, contreparties qui ne répondent pas.
\item Plus d'un prime broker avec des blocages, des limites de ségrégation et de réhypothécation, et le collatéral excédentaire retiré chaque jour.
\item Au moins le déficit le plus important de l'exercice, plus le coût des actions lentes ou incertaines.
\item Une fois par an et après les changements majeurs~; le comité ne doit pas connaître le scénario à l'avance.
\item L'horizon atteint, les décisions et leurs heures, ce qui a cédé, et les correctifs avec leurs responsables et leurs échéances.
\item Six heures sans action~; 1.42~millions de dollars pour l'ensemble le moins coûteux qui atteint 72 heures.
\item Détenir des fonds monétaires et une ligne engagée tirable en quelques heures, répartir le collatéral entre plusieurs courtiers et retirer l'excédent chaque jour~;
lancer l'exercice chaque année et corriger ce qu'il met au jour avant le suivant.
\end{enumerate}
\end{solution}

\begin{solution}{iq:fm:crisis-management:1}
Actionner le kill switch ou l'annulation de masse, arrêter la stratégie, vérifier que les ordres ont disparu, contrôler la position et le P\&L, et appeler le responsable d'incident.
\iqlookfor{d'abord arrêter, puis évaluer.}
\end{solution}

\begin{solution}{iq:fm:crisis-management:2}
Ce qu'elle détient pour notre compte et si c'est ségrégué ou réhypothéqué, nos positions et leur marge ailleurs, où transférer les positions, et notre trésorerie
pour les appels qui suivent.
\iqlookfor{\omterm{def:fm:crisis-management:trapped}{actifs bloqués} et marge de remplacement.}
\end{solution}

\begin{solution}{iq:fm:crisis-management:3}
Actifs liquides vendables le jour même, lignes engagées, transferts de collatéral depuis les comptes en excédent, repo, et réduction des positions, par
ordre de coût par heure de survie gagnée, en vérifiant chaque délai par rapport à l'échéance de l'appel.
\iqlookfor{délais et coûts, pas seulement les montants.}
\end{solution}

\begin{solution}{iq:fm:crisis-management:4}
Le lancer sans prévenir, sur de vraies listes de contacts et d'autorités, avec des défaillances injectées (un signataire absent, un système en panne), et consigner les décisions et
leurs heures par rapport au plan.
\iqlookfor{effet de surprise et réalisme.}
\end{solution}

\begin{solution}{iq:fm:crisis-management:5}
Vendre dans un marché sous tension a un fort impact de marché, et les pertes réalisées peuvent dépasser la marge libérée~; les autres qui vendent les mêmes positions aggravent
la situation.
\iqlookfor{impact et encombrement.}
\end{solution}

\begin{solution}{iq:fm:crisis-management:6}
Choisir un sous-ensemble d'actions dotées de délais, de coûts et d'échéances, pour minimiser le coût sous la contrainte que la trésorerie reste non négative à chaque heure~; une règle gloutonne ignore
les interactions (une action inutile seule mais décisive après d'autres), si bien qu'une petite recherche exhaustive ou en nombres entiers est préférable.
\iqlookfor{contraintes à chaque heure et interactions.}
\end{solution}
